%! AP Statistics — Unit 4, Lesson 4.7 — Group Activity (Answer Key)
\documentclass[10pt]{article}
\usepackage{apstats-article}
\usepackage{apstats-key}

\begin{document}

\pageheader{Unit 4, Lesson 4.7}{Group Activity: Building and Reading Distributions --- Answer Key}
\namepartnerperiod

\vspace{0.1in}

\begin{tcolorbox}[colback=sky!60, colframe=navy, title=\textbf{Tier R --- Remediate},
    arc=2mm, boxrule=0.8pt, left=4mm, right=4mm, top=3mm, bottom=3mm]
\small
\begin{center}
\renewcommand{\arraystretch}{1.4}
\begin{tabular}{@{}cccc@{}}
\textbf{Siblings ($x$)} & \textbf{Frequency} & \textbf{$P(X=x)$} & \textbf{$P(X \le x)$} \\
\hline
0 & 8  & \ans{0.16} & \ans{0.16} \\
1 & 20 & \ans{0.40} & \ans{0.56} \\
2 & 15 & \ans{0.30} & \ans{0.86} \\
3 & 7  & \ans{0.14} & \ans{1.00} \\
\end{tabular}
\end{center}

\begin{enumerate}[leftmargin=1.5em, itemsep=8pt]
  \item Divide each frequency by 50.

  \item \ansline{Valid: all values $\ge 0$ and $0.16+0.40+0.30+0.14=1.00$ \checkmark}

  \item See table above.

  \item \ansline{$P(X \ge 2) = P(2)+P(3) = 0.30+0.14 = 0.44$}

  \item \ansline{0 siblings: $P=0.16$; 3 siblings: $P=0.14$. More likely 0 siblings.}
\end{enumerate}
\end{tcolorbox}

\vspace{0.12in}

\begin{tcolorbox}[colback=goldbg!60, colframe=goldacc, title=\textbf{Tier A --- Approaching Proficiency},
    arc=2mm, boxrule=0.8pt, left=4mm, right=4mm, top=3mm, bottom=3mm]
\small
\begin{center}
\renewcommand{\arraystretch}{1.4}
\begin{tabular}{@{}cccccc@{}}
$x$ & 0 & 1 & 2 & 3 & 4 \\
\hline
$P(X=x)$ & 0.05 & 0.20 & 0.35 & 0.25 & 0.15 \\
\end{tabular}
\end{center}

\begin{enumerate}[leftmargin=1.5em, itemsep=8pt]
  \item \ansline{$0.05+0.20+0.35+0.25+0.15=1.00$; all values $\ge 0$. Valid. \checkmark}

  \item \ansline{$P(1 \le X \le 3)=0.20+0.35+0.25=0.80$. There is an 80\% chance a student reads 1–3 books.}

  \item \ansline{$P(X \ge 3)=0.25+0.15=0.40$. There is a 40\% chance a student reads 3 or more books.}

  \item \ansline{Slightly right-skewed (tail toward 4). Centered around 2 books.}

  \item \ansline{The distribution is slightly right-skewed. Most students read 2 books (35\%); reading 0 is rare (5\%). The values are moderately spread from 0 to 4 books.}
\end{enumerate}
\end{tcolorbox}

\vspace{0.12in}

\begin{tcolorbox}[colback=greenbg!60, colframe=greenacc, title=\textbf{Tier E --- Extension},
    arc=2mm, boxrule=0.8pt, left=4mm, right=4mm, top=3mm, bottom=3mm]
\small
\begin{enumerate}[leftmargin=1.5em, itemsep=8pt]
  \item \ansline{Class A: $0.40+0.35+0.20+0.05=1.00$ \checkmark. Class B: $0.10+0.25+0.40+0.25=1.00$ \checkmark.}

  \item \ansline{Class B tends to participate in more activities. Its probabilities are concentrated at higher values (2 and 3), while Class A is concentrated at 0 and 1.}

  \item \ansline{Class A: $P(X \ge 2)=0.20+0.05=0.25$. Class B: $P(X \ge 2)=0.40+0.25=0.65$. Class B has a much higher probability.}

  \item \ansline{The most common value in Class B is $x=2$. $P(X \ge 2)$ in Class A $= 0.25$.}

  \item \ansline{Class A is right-skewed (most students participate in 0–1 activities). Class B is left-skewed (most students participate in 2–3 activities).}
\end{enumerate}
\end{tcolorbox}

\end{document}
